% Chapitre 4 : écran fictif d'un éditeur avec un fichier Lean et l'infoview.
\begin{figure}[tbp]
\centering
\resizebox{\linewidth}{!}{%
\begin{tikzpicture}[x=1cm,y=1cm,
  num/.style={font=\ttfamily\scriptsize, text=black!35, anchor=east, inner sep=0pt},
  kw/.style={text=tonregle}]
  \fenetre{13.4}{6.4}{Pairs.lean : éditeur}
  % onglet
  \fill[black!4] (0.02,5.35) rectangle (8.2,5.83);
  \fill[white] (0.1,5.35) rectangle (2.1,5.83);
  \draw[black!25] (0.1,5.35) -- (0.1,5.83) -- (2.1,5.83) -- (2.1,5.35);
  \node[font=\sffamily\scriptsize] at (1.1,5.59) {Pairs.lean};
  \draw[black!25] (0.02,5.35) -- (8.2,5.35);
  % le code
  \foreach \n/\y in {1/4.95,2/4.55,3/4.15,4/3.75,5/3.35,6/2.95,7/2.55,8/2.15,9/1.75}
    \node[num] at (0.6,\y) {\n};
  \fill[tonregle!8] (0.7,2.33) rectangle (8.1,2.73);
  \node[mono] at (0.85,4.95) {\textcolor{tonregle}{import} Mathlib};
  \node[mono, text=black!50] at (0.85,4.15) {/-- La somme de deux entiers pairs est paire. -/};
  \node[mono] at (0.85,3.75) {\textcolor{tonregle}{theorem} pair\_add \{a b : ℤ\}};
  \node[mono] at (0.85,3.35) {\ \ \ \ (ha : Even a) (hb : Even b) : Even (a + b) := \textcolor{tonregle}{by}};
  \node[mono] at (0.85,2.95) {\ \ \textcolor{tonregle}{obtain} ⟨k, rfl⟩ := ha};
  \node[mono] at (0.85,2.55) {\ \ \textcolor{tonregle}{obtain} ⟨l, rfl⟩ := hb};
  \draw[black!70, line width=0.6pt] (4.12,2.4) -- (4.12,2.7);
  \node[mono] at (0.85,2.15) {\ \ \textcolor{tonregle}{exact} ⟨k + l, \textcolor{tonregle}{by} ring⟩};
  % barre d'état
  \fill[black!6] (0.02,0.02) rectangle (13.38,0.42);
  \node[font=\sffamily\tiny, text=black!60, anchor=west] at (0.2,0.22)
    {Lean 4 \quad ligne 7, colonne 22 \quad aucune erreur};
  % l'infoview
  \draw[black!30] (8.2,0.42) -- (8.2,5.85);
  \fill[black!3] (8.21,0.43) rectangle (13.37,5.84);
  \node[font=\sffamily\scriptsize\bfseries, text=black!65, anchor=west] at (8.4,5.55) {Infoview};
  \node[font=\sffamily\scriptsize, anchor=west, text=black!60] at (8.4,5.05)
    {▾ \texttt{Pairs.lean\string:7\string:22}};
  \node[font=\sffamily\scriptsize\bfseries, anchor=west] at (8.4,4.55) {Tactic state};
  \node[font=\sffamily\scriptsize, anchor=west, text=black!60] at (8.4,4.15) {1 goal};
  \draw[black!25, fill=white] (8.4,2.7) rectangle (13.15,3.85);
  \node[mono] at (8.55,3.55) {k l : ℤ};
  \node[mono] at (8.55,3.0) {\textcolor{tonregle}{⊢} Even (k + k + (l + l))};
  \node[font=\sffamily\scriptsize\bfseries, anchor=west] at (8.4,2.2) {Messages};
  \node[font=\sffamily\scriptsize, anchor=west, text=black!60] at (8.4,1.8) {aucun};
\end{tikzpicture}}
\caption{Écran fictif : un fichier Lean dans l'éditeur, et l'infoview à droite. Le curseur
est à la fin de la ligne~7 ; l'infoview montre ce qu'il reste à démontrer à cet endroit,
après que les deux hypothèses ont été dépliées.}
\label{fig:ecran-lean}
\end{figure}
